\section{Security Monitoring and Logging (Wazuh SIEM)}

Security monitoring within TAC Industries is designed to provide visibility across infrastructure, authentication events, and system activity without exposing sensitive application data unnecessarily. Centralized monitoring is achieved through a Security Information and Event Management (Wazuh SIEM) platform, which aggregates logs from servers, endpoints, and network devices\cite{wazuh}.

The Wazuh SIEM is used to collect and correlate events such as authentication attempts, privilege escalation activity, configuration changes, and endpoint alerts. By centralizing log ingestion, the platform supports timely detection of anomalous behavior and enables post-event investigation through structured audit trails\cite{nist80092}.

Network level intrusion detection complements endpoint monitoring through Suricata running at the perimeter. Network events are correlated with endpoint and authentication logs within the SIEM, allowing activity to be associated with specific users and systems via VPN identity, directory services, and agent telemetry. This layered visibility supports detection of both endpoint compromise and network-based threats without relying on a single control.

\section{Endpoint and Asset Management (Tactical RMM, Snipe-IT)}

Endpoint and asset management are treated as supporting controls that enable secure operations rather than direct components of the investigative application. Tactical RMM is used to manage enrolled endpoints and servers, providing capabilities such as patch management, system monitoring, and remote maintenance within a controlled administrative scope \cite{tacticalRMM, snipeit}.

Asset ownership and lifecycle tracking are handled through a dedicated asset management platform. This enables hardware and system assets to be recorded, assigned, and audited independently of operational or investigative data. The separation of endpoint management from case data ensures that operational tooling does not become a secondary access path to sensitive information.

Together, these tools support operational resilience, accountability, and system hygiene while maintaining a clear boundary between infrastructure management and investigative workflows.

